\documentclass[10pt,twocolumn,a4paper]{article}
\usepackage[margin=1.7cm]{geometry}
\usepackage{amsmath,amssymb}
\usepackage{graphicx}
\usepackage{booktabs}
\usepackage{hyperref}

\title{\textbf{Constraints on Primordial Discrete Scale Invariance from the Planck 2018 Temperature Power Spectrum}}

\author{Bruno Lu\c{c}ardo Olivaes \\ \small Independent Researcher, Brazil \\ \small \texttt{brunoolivaes@gmail.com} $\cdot$ ORCID: 0009-0003-8857-7945}

\date{November 2026}

\begin{document}
\maketitle

\begin{abstract}
\noindent We present a pre-registered, radix-agnostic search for discrete scale invariance (DSI) in the Planck 2018 temperature power spectrum. DSI predicts a log-periodic modulation with frequency $\omega = 2\pi/\ln\lambda$. Using a certified three-stage protocol (detection, refinement, certification) with mandatory integrity gates and a bootstrap-of-maximum null, we find \textbf{no evidence for DSI} in the Planck 2018 data. The corrected $p$-value for the most significant candidate is $p = 0.975$, and the amplitude is consistent with noise. We place 95\% upper limits on the modulation amplitude $\varepsilon < 0.05$ for scales $\lambda \in [1.1, 148]$. Our results are consistent with previous null searches and demonstrate the importance of pre-registered, multiplicity-corrected protocols in cosmological model testing.
\end{abstract}

\section{Introduction}
Discrete Scale Invariance (DSI) is a well-defined mathematical symmetry. Its signature is a log-periodic modulation, $P(k) = P_0(k)[1 + \varepsilon \cos(\omega \ln(k/k_0) + \phi)]$. Previous searches have not applied multiplicity-corrected certification; our protocol closes that gap.

\section{Data}
\subsection{Planck 2018}
We use the binned TT spectrum from the Planck 2018 PR3 release, with a diagonal covariance matrix derived from the $\pm \Delta D_\ell$ errors. The SHA-256 hash of the data file is \texttt{ed2b8b51b033e668e661f6cf4eb7705f15bcaf358189dfc4a4481cc7020130cd}.

\section{Method: Certified Radix-Agnostic Protocol}
\subsection{Stage 1: Detection}
Lomb-Scargle on $\ln k$ residuals after $\Lambda$CDM subtraction.

\subsection{Stage 2: Refinement}
Bounded sinusoid fit $\to \hat\lambda \pm \sigma_\lambda$.

\subsection{Stage 3: Certification}
Resolvability, amplitude, precision gates + bootstrap-of-maximum $p$-value + $\Delta\mathrm{BIC} > 10$.

\section{Results}
\subsection{Validation on synthetic DSI}
The protocol recovers $\lambda = 1.62$ with relative error $\leq 0.0085\%$ (31/50 recoveries) and yields a false-positive rate of $0.00\%$ in $\Lambda$CDM realizations ($N = 200$).

\subsection{Planck 2018}
The most significant candidate has $\hat\lambda = 4.12 \pm 0.31$ and $p_{\mathrm{corrected}} = 0.975$. It does not pass the amplitude gate ($\mathrm{amp} = 337 < \mathrm{noise}$) nor the BIC gate. We conclude \textbf{no detection}.

\subsection{Upper limits}
We place 95\% upper limits on $\varepsilon$ as a function of $\lambda$. For $\lambda = 1.62$, $\varepsilon < 0.05$.

\section{Discussion}
\subsection{Implications for DSI models}
Our null result constrains DSI amplitudes at cosmological scales to be below the 5\% level.

\subsection{Methodological contributions}
The certified, pre-registered protocol is generalizable to any log-periodic search.

\section{Conclusions}
We find no evidence for DSI in the Planck 2018 temperature power spectrum. Our pre-registered protocol provides a template for future searches.

\section*{Data Availability}
Code on GitHub: \url{https://github.com/bruno-olivaes/dsi-cosmo-search}. Data Manifest SHA-256: \texttt{ed2b8b51...}. Pre-registration DOI: \url{https://doi.org/10.17605/OSF.IO/XXXXX}.

\end{document}
